%% sections/f-constitutional-authority.tex - H. R. 9510 Bill v4.0 (full bill) - App. F
%% Genre: formal Constitutional Authority Statement. Source deliverable 06.
%% Operative sentence in \billbox{}. Single hyphens only; section symbol in tables only.
\billsec{APPENDIX F.}{CONSTITUTIONAL AUTHORITY STATEMENT (NON-OPERATIVE).}\phantomsection\label{toc:aF}

\uscprov{1}{\textit{Independent research draft. Not an enacted law, not pending
legislation, and not legal advice; not endorsed by the FDA, HHS, or any Member of
Congress.}}

\uscprov{1}{Under clause 7(c) of Rule XII of the Rules of the House of
Representatives, a bill or joint resolution may not be introduced unless the sponsor
submits, for printing in the Congressional Record, a statement citing as specifically
as practicable the power or powers granted to Congress in the Constitution to enact
the measure. For H. R. 9510, which amends the Federal Food, Drug, and Cosmetic Act
(21 U.S.C. 301 et seq.), the operative statement is the single sentence below
\cite{house-rule-xii, usc-fdca}.}

\billbox{Congress has the power to enact this legislation pursuant to the following:
Article I, section 8, clause 3, of the Constitution of the United States (the Commerce
Clause), which empowers Congress to regulate commerce among the several States, on
which the Federal Food, Drug, and Cosmetic Act has rested since its enactment in 1938,
as supplemented by Article I, section 8, clause 18 (the Necessary and Proper Clause).}

\uscprov{1}{The paragraphs that follow explain and support that sentence; they are
not part of the statement printed under clause 7(c).}

\uscprov{1}{\textbf{Why the Commerce Clause.} The Federal Food, Drug, and Cosmetic
Act has rested on the Commerce Clause since 1938, and the Act's prohibited-acts
provisions (21 U.S.C. 331) are written expressly in terms of articles introduced
into, or held for sale after shipment in, interstate commerce. Because H. R. 9510
amends that Act, by adding a new section 515D (codified at 21 U.S.C. 360e-5) and
conforming amendments threaded through Title 21, it draws on the same constitutional
source as the statute it modifies. The bill adds no new jurisdictional theory; it
extends a recognized one to a new class of regulated articles \cite{usc-331, usc-321h}.}

\uscprov{1}{\textbf{Devices and their code are articles of interstate commerce.}
Surgical and oncology robots, their components, and the software that governs
robot-patient interaction are manufactured, marketed, shipped, and used across State
lines, and the device authorities the bill builds on, classification under 21 U.S.C.
360c, premarket review under 21 U.S.C. 360e and section 360(k), and predetermined
change control plans under 21 U.S.C. 360e-4, all regulate that interstate market. The
robot-patient interaction code subjected to verification before generation is itself a
traded, updatable article: it is distributed to clinical-investigation sites in
multiple States, modified through interstate software pipelines, and integrated into
devices that move in commerce. Such code and the Physical AI systems that execute it
are both instrumentalities and articles of interstate commerce, squarely within the
reach of clause 3 \cite{usc-360c, usc-360e, usc-360k, usc-360e4}.}

\uscprov{1}{\textbf{Aggregate effects and the Necessary and Proper Clause.} The bill
also reaches activity that, in the aggregate, substantially affects interstate
commerce: Physical AI oncology clinical investigations are conducted by sponsors and
sites operating in a national research market, generating data, evidence, and devices
that enter interstate channels and inform nationwide marketing and Medicare coverage
determinations. Setting a uniform sequencing rule, that automated verification,
validation, and uncertainty quantification must clear before robot-patient code is
generated or executed, is a regulation of the terms on which these articles and
services may be placed in that market. The Necessary and Proper Clause supports the
means chosen: the documentation, attestation, audit-trail, cybersecurity, and
human-oversight duties in section 515D are appropriate and plainly adapted to making
the underlying device and commerce regulation effective.}

\uscprov{1}{\textbf{Filing.} In practice the statement accompanies the measure at the
point of introduction rather than as a separate later filing. The House has
modernized intake through the electronic hopper, and a sponsoring office can generate
and submit the constitutional authority statement, together with the measure, within
that system; the Clerk then makes each statement available to the public in
electronic form. The table below summarizes the requirement \cite{house-ehopper, house-rule-xii}.}

{\footnotesize
\begin{xltabular}{\textwidth}{@{}L{4.2cm} Y@{}}
\hline\noalign{\vskip 0.04cm}
\textbf{Item} & \textbf{Detail}\\
\hline\noalign{\vskip 0.04cm}
\endfirsthead
\hline\noalign{\vskip 0.04cm}
\textbf{Item} & \textbf{Detail}\\
\hline\noalign{\vskip 0.04cm}
\endhead
\hline\noalign{\vskip 0.04cm}
\endlastfoot
Governing rule & Rule XII, clause 7(c), Rules of the House of Representatives\\
\hline\noalign{\vskip 0.04cm}
What is filed & A statement citing as specifically as practicable the constitutional
power or powers to enact the measure\\
\hline\noalign{\vskip 0.04cm}
Primary citation & Article I, \S 8, clause 3 (Commerce Clause)\\
\hline\noalign{\vskip 0.04cm}
Supporting citation & Article I, \S 8, clause 18 (Necessary and Proper Clause)\\
\hline\noalign{\vskip 0.04cm}
Where it appears & Printed in the Congressional Record; made available electronically
by the Clerk\\
\hline\noalign{\vskip 0.04cm}
How it is submitted & With the measure through the electronic hopper at introduction\\
\hline
\end{xltabular}}
\vspace{-0.8cm}
\figcaption{The constitutional authority requirement of Rule XII, clause 7(c): the
primary citation is the Commerce Clause, with the Necessary and Proper Clause in
support, filed with the measure at introduction.}
